% ch-a4-20.tex —— 第 4 篇（Emerton--Gee--Hellmann）续片，底本 PDF 408--419 = 印刷页 395--406
% 本片承接 ch-a4-19（命题 D.1.3 的证明中段），始于附录 D.2，含附录 E、附录 F 与参考文献开头。
% 上一片（ch-a4-19）未闭合命题 D.1.3 的 proof 环境；本片续排证明正文，并在 (4) 证完后 \end{proof}。
\ysrc{395}
由引理 D.1.2 可知，它在 $G$ 中闭且离散，因而在 $X$ 中亦然。这与 $X$ 的紧性矛盾，于是 (3) 得证。

现在设 $H$ 是 $G$ 的子群。按诱导拓扑的定义，$H$ 的开子群都形如 $H\cap U$，其中 $U$ 是 $G$ 的开子群。令 $U_n:=G_n\cap U$，则 $U_n$ 是 $G$ 的开子群，且 $U=\bigcup_{n=1}^\infty U_n$。于是 $H\cap U_n=H_n\cap U_n=H_n\cap U$ 是 $H_n$ 的开子群，而注 D.1.1 表明
\[
\bigcup_{n=1}^\infty H_n\cap U_n=H\cap\bigcup_{n=1}^\infty U_n=H\cap U
\]
也是 $\varinjlim_n H_n$ 的开子群，这就证明了 (4) 的第一个论断。至于拓扑同构 $\varinjlim_n G_n/H_n\isom G/H$，只要比较源与靶的泛映射性质即可立即验证。

显然，若 $H$ 在 $G$ 中闭，则对每个 $n$，$H_n$ 在 $G_n$ 中闭。反之，设对每个 $n$，$H_n$ 在 $G_n$ 中闭。此时每个商 $G_n/H_n$ 仍为 Hausdorff 空间，于是由 (2) 可知 $\varinjlim_n G_n/H_n$ 是 Hausdorff 空间。因此作为典范态射 $G=\varinjlim_n G_n\to\varinjlim_n G_n/H_n$ 的核，$H$ 在 $G$ 中闭，如所断言；这就完成了 (4) 的证明。
\end{proof}

\subsection*{D.2.\quad 可数生成拓扑模}

下面先陈述一个简单的引理。

\begin{lem}{D.2.1}\label{thm:a4-20-1}
若 $A$ 是左 Noether 环（不必交换），则可数生成 $A$-模构成的范畴是所有 $A$-模构成的范畴的阿贝尔子范畴。
\end{lem}

\begin{proof}
唯一不完全清楚的一点是对取子对象封闭。为此注意：若 $N$ 可数生成，则可写 $N=\varinjlim_n N_n$，其中每个 $N_n$ 有限生成。设 $M$ 是 $N$ 的 $A$-子模，则有 $M=\varinjlim_n M_n$，其中 $M_n:=M\cap N_n$。由假设 $A$ 是左 Noether 环，故每个 $M_n$ 仍有限生成，于是 $M$ 确实可数生成。
\end{proof}

现在设 $A$ 是紧 Hausdorff 左 Noether 拓扑环（可能不交换）。此时任意有限生成 $A$-模都带有唯一的紧 Hausdorff 拓扑——称为它的典范拓扑——相对于该拓扑它成为拓扑 $A$-模；而有限生成 $A$-模之间的所有态射关于源与靶上的典范拓扑都连续且严格（即从其定义域到其像诱导拓扑商映射）。（特殊情形下的证明见 [Eme10, Prop. 2.1.3]，该证明立即可推广。要点如下：给定有限表现 $A^{\oplus m}\to A^{\oplus n}\to M\to0$，第一个箭头的像必定闭，因为它是 Hausdorff 空间的紧子集；于是 $M$ 得到自然的紧 Hausdorff 商拓扑，而容易验证它与有限表现的选取无关。用有限表现作类似论证，即得关于态射的论断。）

在我们关心的情形中，$A$ 上的拓扑由开理想组成的 $0$ 的邻域基给出；这时容易看出，任意有限生成 $A$-模上的典范拓扑同样由开 $A$-子模组成的 $0$ 的邻域基给出。在这种情形，我们说 $A$ 与
\ysrc{396}
$M$ 容许线性拓扑，并且从现在起假定 $A$ 上的拓扑确实是线性的。

记 $A^{\oplus\N}$ 为 $A$ 的可数直和。把 $A^{\oplus\N}$ 写成 $A^{\oplus\N}:=\varinjlim_n A^{\oplus n}$（转移映射取显然者），再赋予它归纳极限线性拓扑。

\begin{prop}{D.2.2}\label{thm:a4-20-2}
\begin{enumerate}[label=(\arabic*)]
\item 若 $M$ 是可数生成 $A$-模，则 $M$ 带有唯一拓扑——我们称之为它的典范拓扑——使得某个（等价地，任意）满射 $A^{\oplus\N}\to M$ 是拓扑商映射。赋以典范拓扑后，$M$ 成为线性拓扑 $A$-模。
\item 可数生成 $A$-模之间的态射，关于源与靶上的典范拓扑，必然连续且严格。
\end{enumerate}
\end{prop}

\begin{proof}
若 $M$ 是可数生成 $A$-模，则可写 $M=\varinjlim_n M_n$，其中 $\{M_n\}$ 是 $M$ 的有限生成 $A$-子模的递增序列。若 $\{M'_n\}$ 是 $M$ 的另一个这样的序列，则显然 $\{M_n\}$ 与 $\{M'_n\}$ 相互共尾，即每个 $M_n$ 都包含在某个 $M'_{n'}$ 中，反之亦然。于是，若给每个 $M_n$ 赋以典范拓扑，则由把 $M$ 写成 $M_n$ 的余极限而得到的 $M$ 上的归纳极限线性拓扑与序列 $\{M_n\}$ 的选取无关。我们宣布此拓扑为 $M$ 的典范拓扑。

由于每个 $M_n$ 都以开 $A$-子模组成 $0$ 的邻域基，由注 D.1.1(2) 可知 $M$ 亦然。于是容易推出 $M$ 是拓扑 $A$-模。

若 $f:M\to N$ 是有限生成 $A$-模之间的态射，则可选取 $\{M_n\}$ 与 $\{N_n\}$，使得 $f$ 限制为态射 $M_n\to N_n$。这些态射关于源与靶上的典范拓扑必然连续，因而态射
\[
f:M=\varinjlim_n M_n\longrightarrow\varinjlim_n N_n=N
\]
也连续。

若 $f$ 是单射，则事实上可以先选取 $N_n$，再令 $M_n:=M\cap N_n$。此时 $M_n$ 在 $N_n$ 中闭，于是命题 D.1.3(4) 表明 $M\to N$ 是闭嵌入。

若 $f$ 是满射，则事实上可以先选取 $M_n$，再令 $N_n:=f(M_n)$。同样由命题 D.1.3(4) 可知 $M\to N$ 是拓扑商映射。

总之，我们现在证明了 (2)，而 (1) 是 (2) 的严格性论断的特殊情形。
\end{proof}

\section*{附录 E. $p$-进解析群的表示}

\subsection*{E.1.\quad 光滑 $G$-表示与环 $\OO[[G]]$ 上的模}

设 $G$ 是 $p$-进解析群，并固定 $G$ 的紧开子群 $H$。则通常的完备群环 $\OO[[H]]:=\varprojlim_{J\triangleleft H\ \text{open}}\OO[H/J]$ 是紧（双边）Noether 线性拓扑环，这是 Lazard 的结果 [Laz65]。

\ysrc{397}

\begin{defn}{E.1.1}\label{thm:a4-20-3}
我们记
\[
\OO[[G]]:=\OO[G]\otimes_{\OO[H]}\OO[[H]]=c\text{-}\mathrm{Ind}_H^G\OO[[H]];
\]
先验地，它是 $G$ 的一个（左）表示，兼有相容的 $\OO[[H]]$-右模结构；但按照 [Koh17, §1]（亦见 [Sho20, §3]），它实际带有自然的环结构，该环结构由「$\OO[G]$ 与 $\OO[[H]]$ 皆为子环」这一要求唯一确定。此外，$\OO$-代数 $\OO[[G]]$ 在典范同构意义下与 $H$ 的选取无关。
\end{defn}

任何 $\OO[[G]]$-模都特别是 $\OO[G]$-模，因而可以看作 $\OO$-线性的 $G$-表示。并非每个 $\OO$-线性 $G$-表示都带有与其 $\OO[G]$-作用相容的 $\OO[[G]]$-模结构（这也就是说，并非每个 $\OO[G]$-模都能延拓为 $\OO[[G]]$-模）；但一个重要的事实是：光滑 $G$-表示（下面回忆其定义）确实带有这样的结构。

\begin{defn}{E.1.2}\label{thm:a4-20-4}
\begin{enumerate}[label=(\arabic*)]
\item 若 $M$ 是 $\OO$-线性 $G$-表示，则称元素 $m\in M$ 是光滑的，如果 $m$ 是 $\OO$-挠元，且 $M$ 被 $G$ 的某个开子群固定。\tn{底本此处作「$M$ is fixed by some open subgroup of $G$」，按上下文当系「$m$ is fixed」之误（该条件若施于整个 $M$ 则与 $m$ 无关，且与下一句「由 $m$ 生成的循环模」不协）；译文照录原文（印刷页 397）。}记 $M_{\mathrm{sm}}$ 为 $M$ 的光滑元组成的子集；它是 $M$ 的 $G$-不变 $\OO$-子模。
\item 若 $G$ 在 $\OO$-模 $M$ 上的表示满足 $M_{\mathrm{sm}}=M$，则称它是光滑的。
\end{enumerate}
\end{defn}

若 $M$ 是光滑 $G$-表示，则 $M$ 上的 $\OO[G]$-作用典范地延拓为 $\OO[[G]]$-作用。于是光滑 $\OO$-线性 $G$-表示构成的阿贝尔范畴可以看作 $\OO[[G]]$-模范畴的全子范畴。（见 [Sho20, Lem. 3.5]。）

\begin{rem}{E.1.3}\label{thm:a4-20-5}
按定义，$\OO$-线性 $G$-表示 $M$ 中的一个元素是光滑的，当且仅当把 $M$ 看作 $G$ 的某个（等价地，任意）紧开子群 $H$ 的表示时它是光滑的。又由 $\OO[[H]]$ 上拓扑的定义可知，$m\in M$ 是光滑的，当且仅当由 $m$ 生成的循环 $\OO[[H]]$-模作为 $\OO[[H]]$ 的商是离散的。
\end{rem}

\begin{rem}{E.1.4}\label{thm:a4-20-6}
仅从定义即可看出：若 $M$ 是 $\OO$-线性 $G$-表示，则 $M_{\mathrm{sm}}$ 是 $M$ 的极大光滑子 $G$-表示。此外，若 $M$ 本身就是 $\OO[[G]]$-模，则（用底层的 $\OO[G]$-模结构定义的）$M_{\mathrm{sm}}$ 在赋予其典范 $\OO[[G]]$-模结构后成为 $M$ 的子 $\OO[[G]]$-模。（这可由下述事实推出：若 $H$ 是任意紧 $p$-进解析群，例如 $G$ 的紧开子群，则 $\OO[H]$ 的增广理想生成 $\OO[[H]]$ 的增广理想。）
\end{rem}

\begin{rem}{E.1.5}\label{thm:a4-20-7}
上述大部分讨论（$\OO[[G]]$ 与光滑 $G$-表示的定义，以及把光滑表示实现为一种特殊的 $\OO[[G]]$-模）对任何一个含投射有限开子群的拓扑群 $G$ 都适用。然而我们在下文证明的许多结果（对 $p$-进解析的 $G$）在这种更一般的语境中并不成立。（例 E.2.3 说明了这一点。）
\end{rem}

\ysrc{398}

记 $(\OO[[G]]\text{-}\Mod)^{\mathrm{c.g.}}$ 为 $\OO[[G]]\text{-}\Mod$ 中由可数生成模组成的全子范畴。

\begin{lem}{E.1.6}\label{thm:a4-20-8}
$(\OO[[G]]\text{-}\Mod)^{\mathrm{c.g.}}$ 是 $\OO[[G]]\text{-}\Mod$ 的 Serre 子范畴，且有足够多投射对象。它的投射对象恰是可数生成投射 $\OO[[G]]$-模。
\end{lem}

\begin{proof}
关于 Serre 子范畴的论断中，唯一不平凡的一点是可数生成 $\OO[[G]]$-模的子模仍可数生成。为此注意：对 $G$ 的任意紧开子群 $H$，环 $\OO[[G]]$ 作为 $\OO[[H]]$-模可数生成。于是任意可数生成 $\OO[[G]]$-模作为 $\OO[[H]]$-模也可数生成。由引理 D.2.1，其任意 $\OO[[G]]$-子模作为 $\OO[[H]]$-模可数生成，因而作为 $\OO[[G]]$-模也可数生成。

显然，任意可数生成 $\OO[[G]]$-模都是某个可数生成自由 $\OO[[G]]$-模的商。后者即使在 $\OO[[G]]\text{-}\Mod$ 中也是投射对象，故在 $(\OO[[G]]\text{-}\Mod)^{\mathrm{c.g.}}$ 中投射。这就证明了关于足够多投射对象的论断。于是 $(\OO[[G]]\text{-}\Mod)^{\mathrm{c.g.}}$ 的任意投射对象都是某个可数生成自由模的收缩（retract），从而是可数生成投射模。
\end{proof}

正如在引理 E.1.6 的证明中已经注意到的，任意可数生成 $\OO[[G]]$-模作为 $\OO[[H]]$-模也可数生成，因而按命题 D.2.2 带有典范拓扑（与 $H$ 的选取无关）。

下面记录一个初等引理。

\begin{lem}{E.1.7}\label{thm:a4-20-9}
可数生成拓扑 $\OO[[H]]$-模 $M$ 在其典范拓扑下离散，当且仅当 $M$ 是光滑 $H$-表示。
\end{lem}

\begin{proof}
若 $M$ 离散，且 $m\in M$，则由 $m$ 生成的 $M$ 的 $\OO[[H]]$-子模离散。于是注 E.1.3 表明 $m$ 光滑。故 $M_{\mathrm{sm}}=M$。

反之，设 $M$ 光滑，并选取满射 $\OO[[H]]^{\oplus\N}\to M$。对任意 $n\ge0$，$\OO[[H]]^{\oplus n}$ 的像经过某个离散、从而有限的商，因而具有有限像 $M_n\subseteq M$。我们有 $M=\varinjlim_n M_n$，而按构造，$M$ 上的典范拓扑就是把 $M$ 写成这些 $M_n$ 的余极限时得到的归纳极限线性拓扑。引理 D.1.2 蕴涵 $M$ 离散，如所断言。
\end{proof}

现在定义一个函子，从可数生成 $\OO[[G]]$-模构成的范畴 $(\OO[[G]]\text{-}\Mod)^{\mathrm{c.g.}}$ 到 pro-光滑 $G$-表示构成的范畴 $\mathrm{Pro}\,\mathrm{sm}\,G$，它在对象上由
\[
M\mapsto\mathrm{Pro}(M):=\text{``}\varprojlim\text{''}M/U
\]
给出，其中 $U$ 跑遍 $M$ 的、按 $M$ 上典范拓扑为开的 $\OO[[G]]$-子模所成的余有向集。由引理 E.1.7 可知，$M/U$ 实际上是光滑 $G$-表示。$\mathrm{Pro}(M)$ 的构造具有函子性，这一点来自下述事实：若 $f:M\to N$ 是可数生成 $\OO[[G]]$-模之间的态射，则 $f$ 关于源与靶上的典范拓扑连续。事实上，此时只要 $V$ 是 $N$ 的开 $\OO[[G]]$-子模，$f^{-1}(V)$ 就是 $M$ 的开 $\OO[[G]]$-子模，于是有诱导态射
\[
\mathrm{Pro}(M):=\text{``}\varprojlim\text{''}M/U\to\text{``}\varprojlim\text{''}M/f^{-1}(V)\to\text{``}\varprojlim\text{''}N/V=:\mathrm{Pro}(N).
\]

\ysrc{399}

\begin{prop}{E.1.8}\label{thm:a4-20-10}
\begin{enumerate}[label=(\arabic*)]
\item 函子 $M\mapsto\mathrm{Pro}(M)$ 右正合。
\item 若 $H$ 是 $G$ 的紧开子群，$N$ 是可数生成 $\OO[[H]]$-模，则 $\OO[[G]]\otimes_{\OO[[H]]}N$ 对左导出函子 $L^i\mathrm{Pro}$ 无环。
\item 所有可数生成光滑表示都对左导出函子 $L^i\mathrm{Pro}$ 无环。
\end{enumerate}
\end{prop}

\begin{proof}
设 $0\to M'\to M\to M''\to0$ 是可数生成 $\OO[[G]]$-模的短正合列。若 $U$ 是 $M$ 的开 $\OO[[G]]$-子模，令 $U'=M'\cap U$，并以 $U''$ 记 $U$ 在 $M''$ 中的像。由于 $M\to M''$ 是开的 $G$-等变态射，当 $U$ 跑遍 $M$ 的所有 $G$-不变开子模时，诸像 $U''$ 跑遍 $M''$ 的所有 $G$-不变开子模。于是得到短正合列
\[
0\to\text{``}\varprojlim\text{''}M'/U'\to\mathrm{Pro}(M)\to\mathrm{Pro}(M'')\to0.
\]
存在显然的满射 $\mathrm{Pro}(M')\to\text{``}\varprojlim\text{''}M'/U'$，因此
\begin{equation}\tag{E.1.9}\label{eq:a4-20-1}
\mathrm{Pro}(M')\to\mathrm{Pro}(M)\to\mathrm{Pro}(M'')\to0
\end{equation}
正合，如所断言。

现在证明 (2)。设 $N$ 是可数生成 $\OO[[H]]$-模，则可找到由可数生成自由 $\OO[[H]]$-模组成的消解 $P^\bullet\to N$；于是 $\OO[[G]]\otimes_{\OO[[H]]}P^\bullet$ 是由可数生成自由 $\OO[[G]]$-模给出的 $M:=\OO[[G]]\otimes_{\OO[[H]]}N$ 的消解。\tn{底本此处作「countable generated free $\OO[[G]]$-modules」，当为「countably generated」之误；译文照录原文（印刷页 399）。}（回忆 $\OO[[G]]$ 在 $\OO[[H]]$ 上平坦。）用这个消解计算左导出函子 $L^i\mathrm{Pro}$，再应用下文的引理 E.1.10(2)，即得 (2) 的论断。

最后证明 (3)。先设 $M''$ 是可数生成光滑 $G$-表示，并设 $M$ 是可数生成自由 $\OO[[G]]$-模，带有满射 $M\to M''$。引理 E.1.7 表明该满射的核 $M'$ 在 $M$ 中开，于是下文的引理 E.1.10(1) 表明 $L^1\mathrm{Pro}(M'')=0$；因此 $L^1\mathrm{Pro}$ 在可数生成光滑表示上消没。

现在设 $H$ 是 $G$ 的紧开子群。我们可以找到某个光滑 $H$-表示 $N$，以及满射 $M:=\OO[[G]]\otimes_{\OO[[H]]}N\to M''$。\tn{底本此处作「We may write find a surjection」，衍一动词；译文按「我们可以找到满射」理解（印刷页 399）。}该满射的核 $N'$ 仍可数生成且光滑；又因 (2) 表明 $M$ 是 $L^\bullet\mathrm{Pro}$-无环的，我们得到 $i\ge1$ 时的同构 $L^{i+1}\mathrm{Pro}(M'')\isom L^i\mathrm{Pro}(M')$。再作显而易见的维数平移论证，并计入上一段的结果，即得 (3)。
\end{proof}

我们不清楚 (E.1.9) 在一般情形下是否左正合；这样的左正合性等价于诸交 $U'$ 在 $M'$ 的所有 $G$-不变开子模中构成共尾族。（若去掉 $G$-不变性条件，这一点总成立，因为 $M'\to M$ 必是拓扑嵌入。）下述引理给出了一些它确实成立的场合。

\ysrc{400}

\begin{lem}{E.1.10}\label{thm:a4-20-11}
若下列两个条件之一成立，则 (E.1.9) 左正合：
\begin{enumerate}[label=(\arabic*)]
\item $M'$ 在 $M$ 中开；或
\item 嵌入 $M'\into M$ 形如
\[
M':=\OO[[G]]\otimes_{\OO[[H]]}N'\into\OO[[G]]\otimes_{\OO[[H]]}N=:M,
\]
其中 $N'\into N$ 是 $\OO[[H]]$-模的嵌入（对某个紧开子群 $H\le G$）。
\end{enumerate}
\end{lem}

\begin{proof}
(1) 若 $M'$ 在 $M$ 中开，则 $M'$ 的任意开 $G$-不变子模 $U'$ 本身也是 $M$ 的开 $G$-不变子模，于是可取 $U=U'$。因此在这种情形 (E.1.9) 正合。

(2) 若 $V'$ 是 $M'$ 的开 $G$-不变子模，则 $W':=N'\cap V'$ 是 $N'$ 的开 $\OO[[H]]$-子模，于是 $U':=\OO[[G]]\otimes_{\OO[[H]]}W'$ 是 $M'$ 的、含于 $V'$ 的开 $G$-不变子模；因此这种形式的子模在 $M'$ 的所有开 $G$-不变子模中构成共尾族。又因 $N'\into N$ 必是拓扑嵌入，可以找到 $N$ 的开 $\OO[[H]]$-子模 $W$，使得 $W'=N'\cap W$。于是 $U:=\OO[[G]]\otimes_{\OO[[H]]}W$ 是 $M$ 的开 $G$-不变子模，且满足 $M'\cap U=U'$。因此在这种情形 (E.1.9) 也正合。
\end{proof}

\subsection*{E.2.\quad 光滑表示的导出范畴}

设 $G$ 是 $p$-进解析群，$\OO[[G]]$ 是定义 E.1.1 中的环。以 $D(\OO[[G]])$ 记与 $\OO[[G]]$-模范畴相联系的稳定 $\infty$-范畴。它带有标准的 $t$-结构，关于此结构它既左完备又右完备。（正如我们在 §A.5 中指出的，对任何环上模范畴所联系的稳定 $\infty$-范畴，这一点都成立。）

\begin{defn}{E.2.1}\label{thm:a4-20-12}
以 $D_{\mathrm{sm}}(\OO[[G]])$ 记 $D(\OO[[G]])$ 的全子范畴，其对象在每一阶的上同调都是光滑 $G$-表示。
\end{defn}

$D(\OO[[G]])$ 上的 $t$-结构诱导 $D_{\mathrm{sm}}(\OO[[G]])$ 上的 $t$-结构，其心是光滑 $G$-表示构成的阿贝尔范畴 $\mathrm{sm}\,G$。$\OO[[G]]\text{-}\Mod$ 的全子范畴 $\mathrm{sm}\,G$ 对取子对象、商以及 $\OO[[G]]$ 中的余极限封闭。因此 $\mathrm{sm}\,G$ 仍是 Grothendieck 阿贝尔范畴，特别是它有足够多内射对象。

\begin{prop}{E.2.2}\label{thm:a4-20-13}
存在典范等价 $D(\mathrm{sm}\,G)\isom D_{\mathrm{sm}}(\OO[[G]])$。
\end{prop}

\begin{proof}
如上所述，把光滑 $G$-表示作为 $\OO[[G]]$-模的全子范畴嵌入，就把前者实现为后者的局部化子范畴。以 $F$ 记此函子，并先证明 $F$ 满足命题 A.5.7 的等价条件。为此只需证明：若 $X$ 是光滑 $G$-表示，$Y$ 是内射光滑 $G$-表示，则可找到光滑 $G$-表示 $Z$ 及满射 $Z\to X$，使得 $i>0$ 时 $\mathrm{Ext}^i_{\OO[[G]]}(Z,Y)=0$。

为此选取 $G$ 的紧开子群 $H$，并令 $Z=c\text{-}\mathrm{Ind}_H^GX=\OO[[G]]\otimes_{\OO[[H]]}X$，它自然满射到 $X$。由于忘函子 $\OO[[G]]\text{-}\Mod\to\OO[[H]]\text{-}\Mod$ 有正合左伴随，即 $\OO[[G]]\otimes_{\OO[[H]]}-$，它保持内射对象。同一个函子限制到 $\mathrm{sm}\,H$ 上时可以解释为 $c\text{-}\mathrm{Ind}_H^G$，从而给出忘函子 $\mathrm{sm}\,G\to\mathrm{sm}\,H$ 的左伴随；于是
\ysrc{401}
后一函子也保持内射对象。因此当 $Y$ 是 $\mathrm{sm}\,H$ 中的内射对象时，$\mathrm{Ext}^i_{\OO[[G]]}(c\text{-}\mathrm{Ind}_H^GX,Y)=\mathrm{Ext}^i_{\OO[[H]]}(X,Y)$；于是问题归结为验证：若 $X,Y$ 是 $\mathrm{sm}\,H$ 的对象且 $Y$ 内射，则 $i>0$ 时 $\mathrm{Ext}^i_{\OO[[H]]}(X,Y)=0$。简言之，我们把原问题中的 $G$ 换成了它的紧开子群 $H$。在这种情形，所需论断可由推论 B.1.16 推出（把它用在例 B.1.1 的语境中）。

下面注意：嵌入 $\mathrm{sm}\,G\into\OO[[G]]\text{-}\Mod$ 有右伴随（由伴随函子定理它必存在，因为光滑表示的余极限仍光滑），而这个右伴随有相当显式的描述，即
\[
M\mapsto M_{\mathrm{sm}}:=\operatorname*{colim}_{H,n}M[\varpi^n]^H,
\]
其中在余极限中，$n$ 跑遍正整数，$H$ 跑遍 $G$ 的开子群。由命题 A.7.1，我们于是可以构造相应的右导出函子 $R\Gamma_{\mathrm{sm}}:D(\OO[[G]])\to D(\mathrm{sm}\,G)$，它是 $F$ 的右伴随。此外，$R\Gamma_{\mathrm{sm}}$ 具有有限的上同调振幅。\footnote{若 $G$ 的维数为 $d$，且 $H$ 充分小（即它是一致 pro-$p$ 群），则 $M\mapsto M[\pi^n]^H$ 的导出函子具有上同调振幅 $d+1$，$R\Gamma_{\mathrm{sm}}$ 亦然（它是这些不同导出函子的滤过余极限）。（译者注：本脚注有两处与正文不一致：其一，$M[\pi^n]^H$ 中作 $\pi$，而正文同一定义处（印刷页 401）作 $\varpi$；其二，「as cohomological amplitude」中的 as 当为 has 之误。译文照录原文（印刷页 401 脚注 34）。）}于是命题由命题 A.7.3 推出。
\end{proof}

\begin{exam}{E.2.3}\label{thm:a4-20-14}
下述（反）例将说明，命题 E.2.2 确实依赖于 $G$ 是 $p$-进解析群这一假设，或至少依赖于完备群环 $\OO[[G]]$ 作为 pro-Artin 环的若干性质，而这些性质对更一般的局部 pro-$p$ 群未必成立。

作为例子，取 $G=\prod_{i=0}^\infty\Z/p\Z$，带其自然的投射有限拓扑。为便于与所引结果作比较，我们改用系数 $k$ 而非 $\OO$；但读者容易说服自己，这一系数改变不起本质作用，且上述讨论容易推广到 $\OO$-系数的情形。乘积 $G_n:=\prod_{i=0}^n\Z/p\Z$ 的群环有下述描述：$k[G_n]=k[x_0,\dots,x_n]/(x_0^p,\dots,x_n^p)$；因此完备群环 $k[[G]]$ 有描述
\[
\begin{aligned}
k[[G]]&=k[[x_0,\dots,x_n,\dots]]/(x_0^p,\dots,x_n^p,\dots)\\
&:=\varprojlim_n k[[x_0,\dots,x_n]]/(x_0^p,\dots,x_n^p).
\end{aligned}
\]
若以 $\mathfrak m$ 记 $k[[G]]$ 的极大理想（即增广理想），则尽管 $\mathfrak m$ 在 $k[[G]]$ 中开，$k[[G]]$ 上的自然投射有限拓扑仍弱于它的 $\mathfrak m$-进拓扑。例如 $\mathfrak m/\mathfrak m^2=\varprojlim_n k\langle x_0,\dots,x_n\rangle$ 是无限维的，而 $k[[G]]$ 上的投射有限拓扑赋予 $\mathfrak m/\mathfrak m^2$ 的是其自然投射有限拓扑，而非离散拓扑。

以 $R$ 记 $k[[G]]$ 的多项式子环 $k[x_0,\dots,x_n,\dots]$。光滑 $G$-表示构成的范畴可以等同于这样的 $R$-模范畴：其中每个元素都被某个理想 $(x_i)_{i\ge n}$（对某个 $n\ge0$）零化。这是 [Nee11, Construction 1.1] 中所引入的、记作 $\mathcal A$ 的那类范畴的一个例子；而 [Nee11, Thm. 1.1] 随后表明 $D(\mathrm{sm}\,G)$ 不左完备，[Nee11, Rem. 1.2] 则表明 $D(\mathrm{sm}\,G)$ 中可数积的构造具有无界上同调振幅（这就说明了命题 A.5.4 中某个有界振幅判据的必要性）。因此，命题 E.2.2 的证明中那个把结果从 $D^+(\mathrm{sm}\,G)$ 推广到 $D(\mathrm{sm}\,G)$ 的论证失效。
\end{exam}

\ysrc{402}

事实上，就连典范函子 $F:D^+(\mathrm{sm}\,G)\to D^+(k[[G]])$ 也不全忠实。确实，$G$ 的一维平凡表示对应于 $k[[G]]$-模 $k:=k[[G]]/\mathfrak m$，于是 $\mathrm{Ext}^1_{k[[G]]}(k,k)=\mathrm{Hom}_k(\mathfrak m/\mathfrak m^2,k)$（到 $k$ 的 $k$-线性对偶）。另一方面，$\mathrm{Ext}^1_{\mathrm{sm}\,G}(k,k)$ 对应于 $\mathrm{Hom}_k(\mathfrak m/\mathfrak m^2,k)$ 中由这样的泛函构成的真子空间：当 $\mathfrak m/\mathfrak m^2$ 带有其自然投射有限拓扑时，这些泛函连续。

\subsection*{E.3.\quad 紧对象与凝聚对象}

我们先给出一些预备知识，依据 [Sho20, §§2,3] 的结果。

\begin{lem}{E.3.1}\label{thm:a4-20-15}
设 $M$ 是光滑 $\OO$-线性 $G$-表示，则下列条件等价：
\begin{enumerate}[label=(\arabic*)]
\item $M$ 作为 $\OO[G]$-模有限生成。
\item $M$ 作为 $\OO[[G]]$-模有限生成。
\item 存在 $G$ 的紧开子群 $H$、$H$ 在某个有限生成挠 $\OO$-模 $N$ 上的光滑表示，以及满射 $c\text{-}\mathrm{Ind}_H^GN\to M$。
\item 对 $G$ 的每个紧开子群 $H$，都存在 $H$ 在某个有限生成挠 $\OO$-模 $N$ 上的光滑表示，以及满射 $c\text{-}\mathrm{Ind}_H^GN\to M$。
\end{enumerate}
\end{lem}

\begin{proof}
显然 (4) 蕴涵 (3)，(3) 蕴涵 (1)，(1) 蕴涵 (2)。于是设 (2) 成立，并设 $S\subset M$ 是 $M$ 作为 $\OO[[G]]$-模的一个有限生成集。若 $H$ 是 $G$ 的任意紧开子群，$N$ 是由 $S$ 生成的 $M$ 的 $\OO[[H]]$-子模（它与由 $S$ 生成的 $\OO[H]$-子模相同，因为 $M$ 的光滑性蕴涵 $H$ 的某个开子群在 $S$ 上平凡作用），则 $c\text{-}\mathrm{Ind}_H^G=\OO[G]\otimes_{\OO[H]}N$ 满射到 $M$，这就证明了 (4)。\tn{底本此处作「$c\text{-}\mathrm{Ind}_H^G=\OO[G]\otimes_{\OO[H]}N$」，既漏写了 $c\text{-}\mathrm{Ind}_H^G$ 的上标 $N$，张量积下标亦作 $\OO[H]$（按上下文当为 $\OO[[H]]$）；译文照录原文（印刷页 402）。}
\end{proof}

\begin{defn}{E.3.2}\label{thm:a4-20-16}
若光滑 $G$-表示满足引理 E.3.1 的等价条件，则称它有限生成。
\end{defn}

\begin{defn}{E.3.3}\label{thm:a4-20-17}
设 $M$ 是 $\OO$-模。若存在光滑 $G$-表示的短正合列
\[
M_1\to M_2\to M\to0,
\]
其中 $M_1,M_2$ 有限生成，则称 $G$ 在 $M$ 上的光滑表示有限表现。
\end{defn}

下述结果是 [Sho20, Prop. 3.8]。

\begin{lem}{E.3.4}\label{thm:a4-20-18}
光滑 $G$-表示有限表现（按定义 E.3.3 的意义），当且仅当它作为 $\OO[[G]]$-模有限表现。
\end{lem}

现在我们得到下述结果。

\begin{prop}{E.3.5}\label{thm:a4-20-19}
范畴 $\mathrm{sm}\,G$ 紧生成。$\mathrm{sm}\,G$ 的紧对象恰是有限表现的光滑表示。
\end{prop}

\begin{proof}
我们概述证明，这是标准论证。首先注意：若 $M$ 是光滑 $\OO$-线性 $G$-表示，$\{M_i\}$ 是 $M$ 的所有有限生成子表示构成的（有向！）集合，则显然 $M=\bigcup_i M_i=\operatorname*{colim}M_i$。对每个 $i$，选取满射 $c\text{-}\mathrm{Ind}_{H_i}^GN_i\to M_i$，其中 $H_i$ 在 $G$ 中紧开，$N_i$ 是 $H_i$ 在某个有限生成挠 $\OO$-模上的光滑表示。若 $M'_i$ 记该满射的核，则如我们刚才所见，$M'_i=\operatorname*{colim}(M'_i)_j$，其中 $(M'_i)_j$ 跑遍 $M'_i$ 的有限生成子表示。

\ysrc{403}

于是 $M=\operatorname*{colim}_{i,j}(c\text{-}\mathrm{Ind}_{H_i}^GN_i)/(M'_i)_j$，这就把 $M$ 表为有限表现表示的余极限。因此，只要证明有限表现表示是紧对象，便证明了 $\mathrm{sm}\,G$ 紧生成。

现在设 $M$ 是紧对象。把 $M$ 表为有限表现表示的余极限并应用其紧性，可知 $M$ 是某个有限表现表示的直和项，因而容易看出它本身也是有限表现表示（可以直接论证，也可以借助引理 E.3.4）。反之，若 $M$ 有限表现，则引理 E.3.4 表明 $M$ 是紧对象（甚至在 $\OO[[G]]$-模范畴中也是）。
\end{proof}

我们一般不知道有限表现的光滑 $G$-表示是否构成 $\mathrm{sm}\,G$ 的阿贝尔子范畴，尽管看起来很可能如此；参见猜想 6.1.4，以及注 6.1.6 中对这一方向上已知结果的讨论。如果确实如此，则 $\mathrm{sm}\,G$ 是局部凝聚阿贝尔范畴，从而 $D(\mathrm{sm}\,G)$ 上的 $t$-结构凝聚。这时（用引理 A.6.5）可知，$D(\mathrm{sm}\,G)$ 中凝聚对象构成的子范畴恰是（上同调）有界、且各阶上同调有限表现的复形构成的全子范畴 $D^b_{\mathrm{f.p.}}(\mathrm{sm}\,G)$。

\subsection*{E.4.\quad 作为 pro-光滑复形的 $\OO[[G]]$-模复形}

形式的 pro-范畴 $\mathrm{Pro}\,\mathrm{sm}\,G$ 有足够多投射对象。于是可以构造与之相联系的稳定 $\infty$-范畴 $D^-(\mathrm{Pro}\,\mathrm{sm}\,G)$，它带自然的 $t$-结构。

我们也可以构造形式的 pro-$\infty$-范畴 $\mathrm{Pro}\,D^b(\mathrm{sm}\,G)$。\tn{底本此处作「We may also from the formal Pro-$\infty$-category」，当为「We may also form the formal Pro-$\infty$-category」之误；译文照录原文（印刷页 403）。}$D^b(\mathrm{sm}\,G)$ 上的 $t$-结构诱导 $\mathrm{Pro}\,D^b(\mathrm{sm}\,G)$ 上的左完备 $t$-结构（这是命题 A.6.12 的对偶陈述），其心等价于 $\mathrm{Pro}\,\mathrm{sm}\,G$。于是（由命题 A.5.7 的对偶版本，即 [Lur17, Prop. 1.3.3.7]）存在典范函子
\begin{equation}\tag{E.4.1}\label{eq:a4-20-2}
D^-(\mathrm{Pro}\,\mathrm{sm}\,G)\to\mathrm{Pro}\,D^b(\mathrm{sm}\,G).
\end{equation}

回忆 $(\OO[[G]]\text{-}\Mod)^{\mathrm{c.g.}}$ 记可数生成 $\OO[[G]]$-模构成的阿贝尔范畴。由引理 E.1.6，这是有足够多投射对象（即诸可数生成自由 $\OO[[G]]$-模）的阿贝尔范畴，因此我们可以考虑它的有界上方复形构成的稳定 $\infty$-范畴 $D^-\bigl((\OO[[G]]\text{-}\Mod)^{\mathrm{c.g.}}\bigr)$。又因（由引理 E.1.6）$(\OO[[G]]\text{-}\Mod)^{\mathrm{c.g.}}$ 事实上是 $\OO[[G]]\text{-}\Mod$ 的 Serre 子范畴，考虑 $D^-_{\mathrm{c.g.}}(\OO[[G]])$ 也有意义，它是 $D^-(\OO[[G]])$ 的、由各阶上同调可数生成的复形构成的全子稳定 $\infty$-范畴。$D^-(\OO[[G]])$ 上的典范 $t$-结构诱导 $D^-_{\mathrm{c.g.}}(\OO[[G]])$ 上的 $t$-结构，其心恰是 $(\OO[[G]]\text{-}\Mod)^{\mathrm{c.g.}}$。

\begin{lem}{E.4.2}\label{thm:a4-20-20}
存在典范等价
\begin{equation}\tag{E.4.3}\label{eq:a4-20-3}
D^-\bigl((\OO[[G]]\text{-}\Mod)^{\mathrm{c.g.}}\bigr)\isom D^-_{\mathrm{c.g.}}(\OO[[G]]).
\end{equation}
\end{lem}

\begin{proof}
这可由命题 A.5.7 的对偶版本，即 [Lur17, Prop. 1.3.3.7] 推出：$(\OO[[G]]\text{-}\Mod)^{\mathrm{c.g.}}$ 中的投射对象恰是可数生成投射 $\OO[[G]]$-模，而它们可用来在 $\OO[[G]]\text{-}\Mod$ 中计算 Ext。
\end{proof}

现在对命题 E.1.8 中的函子作左导出，得到函子
\[
L\mathrm{Pro}:D^-\bigl((\OO[[G]]\text{-}\Mod)^{\mathrm{c.g.}}\bigr)\to D^-(\mathrm{Pro}\,\mathrm{sm}\,G).
\]

\ysrc{404}

先用 (E.4.3) 中等价的一个逆作前复合，再用 (E.4.1) 作后复合，便得到函子
\begin{equation}\tag{E.4.4}\label{eq:a4-20-4}
D^-_{\mathrm{c.g.}}(\OO[[G]])\to\mathrm{Pro}\,D^b(\mathrm{sm}\,G).
\end{equation}
以 $D^b_{\mathrm{c.g.}}(\mathrm{sm}\,G)$ 记 $D^b(\mathrm{sm}\,G)$ 中由各阶上同调可数生成的对象构成的全子范畴，并以 $D^b_{\mathrm{c.g.},\mathrm{sm}}(\OO[[G]])$ 记 $D^b(\OO[[G]])$ 中由各阶上同调可数生成且光滑的复形构成的全子范畴。命题 E.2.2 的等价显然诱导出等价
\begin{equation}\tag{E.4.5}\label{eq:a4-20-5}
D^b_{\mathrm{c.g.}}(\mathrm{sm}\,G)\isom D^b_{\mathrm{c.g.},\mathrm{sm}}(\OO[[G]]).
\end{equation}
由命题 E.1.8(3) 可知，等价 (E.4.5) 与（函子 (E.4.4) 在 $D^b_{\mathrm{c.g.},\mathrm{sm}}(\OO[[G]])$ 上的限制）的复合，等价于典范函子
\[
D^b_{\mathrm{c.g.}}(\mathrm{sm}\,G)\into D^b(\mathrm{sm}\,G)\into\mathrm{Pro}\,D^b(\mathrm{sm}\,G)
\]
的复合。

总之，我们有 (E.4.4) 中的函子 $D^-_{\mathrm{c.g.}}(\OO[[G]])\to\mathrm{Pro}\,D^b(\mathrm{sm}\,G)$；把它限制到 $D_{\mathrm{c.g.}}(\mathrm{sm}\,G)^b$（借助等价 (E.4.5) 视为 $D^-_{\mathrm{c.g.}}(\OO[[G]])$ 的子范畴）上，它等价于显然的包含函子。

\section*{附录 F. 群表示的导出模叠}

本附录简要介绍伽罗瓦表示的导出模叠，并说明其动机。细节参见 [Zhu25a, §2]。

令
\[
\Gamma:=\langle g_1,\dots,g_m\mid r_1=\dots=r_n=1\rangle
\]
是有限表现群，生成元为 $g_1,\dots,g_m$，关系为 $r_1,\dots,r_n$。（应当记住的动机性例子是局部域或整体域的有限扩张的伽罗瓦群；在正文中我们转而考虑各种拓扑有限表现的投射有限群，例如由 $p$-进局部域的绝对伽罗瓦群给出的群，但为了说明动机，字面意义上的有限群这一情形已足够。）

于是，给出表示 $\Gamma\to\mathrm{GL}_d(A)$（对某个环 $A$）相当于给出满足诸关系 $r_j$ 的元素 $x_1,\dots,x_m\in\mathrm{GL}_d(A)$。以 $V_\Gamma$ 记 $(\mathrm{GL}_d)^m$ 中分类满足诸关系 $r_j$ 的 $m$ 元组的闭子概形。$\mathrm{GL}_d$ 通过同时共轭作用在 $V_\Gamma$ 上，商叠 $\XX_\Gamma:=[V_\Gamma/\mathrm{GL}_d]$ 因而是 $\Gamma$ 的 $d$ 维表示的模叠（这里写在任意交换环上，尽管我们当然也可以在任意给定的基环 $\OO$ 上工作，然后考虑 $\OO$-代数上的表示）。它是代数叠（在 $\Z$ 上，或在所选的任何其他基环 $\OO$ 上）。

若固定一个具体的表示 $\rho:\Gamma\to\mathrm{GL}_d(l)$（定义在某个域 $l$ 上），则该表示给出 $\XX_\Gamma$ 的一个 $l$-值点，而 $\rho$ 的提升环 $R_\rho$ 是 $\XX_\Gamma$ 在点 $\rho$ 处的通用（versal）环。$\XX_\Gamma$ 在 $\rho$ 处的切空间可以用群上同调描述为 $H^1(\Gamma,\mathrm{Ad}\rho)$，而 $H^0(\Gamma,\mathrm{Ad}\rho)$ 计算 $\rho$ 的无穷小自同构。转到更高次上同调，已知 $H^2(\Gamma,\mathrm{Ad}\rho)$ 的元素给出 $\rho$ 的高阶形变的阻碍。然而，这套阻碍理论是额外的一份数据——它
\ysrc{405}
并不由叠 $\XX_\Gamma$ 本身内蕴地决定。例如，一般说来人们无法用 $\XX_\Gamma$ 来「计算」$H^2(\Gamma,\mathrm{Ad}\rho)$（这与 $H^0$、$H^1$ 不同）。

当我们用群上同调来描述提升环 $R_\rho$ 时，这一点会具体地显现出来：若考虑 $l$-代数上的提升，并记 $h^i:=\dim_lH^i(\Gamma,\mathrm{Ad}\rho)$，则可以把 $R_\rho$ 写成商 $l[[x_1,\dots,x_{h^1-h_0+d^2}]]/(f_1,\dots,f_{h^2})$。\tn{底本此处作 $h^1-h_0+d^2$，其中 $h_0$ 排成下标；按定义 $h^i:=\dim_lH^i(\Gamma,\mathrm{Ad}\rho)$，此处当为 $h^0$。译文照录原文（印刷页 405）。}当 $\XX_\Gamma$ 在 $\rho$ 处具有「预期维数」$-h^0+h^1-h^2$ 时，诸 $f_i$ 必在 $l[[x_1,\dots,x_{h^1-h^0+d^2}]]$ 中构成正则序列，且 $R_\rho$ 是局部完全交。但若 $\XX_\Gamma$ 的维数大于预期，则 $R_\rho$ 的环论结构就不那么清楚，它与量 $h^2$ 的联系也不那么清楚。

因此，考虑 $\XX_\Gamma$ 的导出类比的一个理由是：它会收紧群上同调与 $\XX_\Gamma$ 在诸点 $\rho$ 处的无穷小结构之间的联系。导出表示叠的精确定义相当复杂（见 [Zhu25a, §2]），但大致如下：正如导出情形中通常的做法，我们并不只计入 $\Gamma$ 中出现的诸关系 $r_i$，而是考虑「更高阶的关系」，即关系之间的关系，如此等等。在群未必交换的非阿贝尔情形，这一点可以形式地表述为：找到 $\Gamma$ 的、由自由群（或者，出于技术原因，实际上由自由幺半群）给出的单纯消解 $\Gamma_\bullet$。分类同态 $\Gamma_n\to\mathrm{GL}_d(A)$（系数环为 $A$）的概形 $V_{\Gamma_n}$ 于是形如 $\operatorname{Spec}R_{\Gamma_n}$，其中诸环 $R_{\Gamma_n}$ 是多项式环的局部化（因为诸 $\Gamma_n$ 是自由幺半群），而这些环本身又组织成一个单纯交换环 $R_{\Gamma_\bullet}$。我们可以考虑相应的导出仿射概形 $V_{\Gamma_\bullet}$，并构造导出代数叠 $\XX_{\Gamma_\bullet}:=[V_{\Gamma_\bullet}/\mathrm{GL}_d]$，其底层的经典子叠与 $\XX_\Gamma$ 一致。在合适的 $\infty$-范畴框架中（例如见 [Zhu25a, §2] 的讨论），导出叠 $\XX_{\Gamma_\bullet}$ 的定义与消解 $\Gamma_\bullet$ 的选取无关。

这里有一个标准而有用的类比：把从经典环到导出环的过渡（或用我们采用的 [ČS24] 的术语，「动态环」），比作从约化环到任意环的过渡，后者出现在从簇到概形的过渡中；而把从经典概形或叠到导出概形或叠的过渡，比作从簇到未必约化的概形的过渡。即使对域 $l$ 上的经典簇，考虑它在非约化环上的点也很有意思：例如其切空间可通过 $l[\epsilon]/(\epsilon^2)$-值点来计算。类似地，对导出叠 $\XX$（即使它实际上是经典的），也存在 $l[\epsilon]/(\epsilon^2)$ 的非经典类比物，它们在经典或导出叠 $\XX$ 中的点决定了 $\XX$ 的阻碍理论；利用这样的点，我们可以把切空间的经典计算推广为对 $\XX$ 在点 $x$ 处的余切复形 $\mathbf L_{\XX,x}$ 的计算。特别地，当 $\XX=\XX_{\Gamma_\bullet}$ 时，可得
\[
\mathbf L_{\XX_{\Gamma_\bullet},\rho}=C^\bullet(\Gamma,\mathrm{Ad}\rho)^*[-1]
\]
（这是计算群上同调 $H^\bullet(\Gamma,\mathrm{Ad}\rho)$ 的上链复形的对偶的一个平移）——例如见 [Zhu25a, §2.2]。这把 $H^0$ 与 $H^1$ 在经典情形中的几何解释推广开来。

当 $i>2$ 时 $H^i(\Gamma,\mathrm{Ad}\rho)=0$ 的情形下，这尤其有用：此时 $\mathbf L_{\XX_{\Gamma_\bullet},\rho}$ 的上同调只出现在次数 $[-1,1]$ 中。这意味着 $\XX_{\Gamma_\bullet}$ 在 $\rho$ 的邻域内拟光滑（即导出局部完全交）。具体地说，这意味着存在形式光滑态射
\[
\operatorname{Spec}l[[x_1,\dots,x_{h^1-h^0+d^2}]]/(f_1,\dots,f_{h^2})\to\XX_{\Gamma_\bullet},
\]

\ysrc{406}

它把定义域的闭点送到 $\rho$，只要把定义域按导出的意义理解：若诸 $f_i$ 不构成正则序列，则环 $l[[x_1,\dots,x_{h^1-h_0+d^2}]]/(f_1,\dots,f_{h^2})$ 将是非经典的动态环（其高阶同伦群由序列 $(f_i)$ 的 Koszul 复形计算）。于是我们看到，$\XX_{\Gamma_\bullet}$ 在 $\rho$ 的邻域内实际上是经典的，当且仅当底层的经典叠 $\XX_\Gamma$ 在这样的邻域内具有预期维数 $-h_0+h_1-h_2=\chi(\mathbf L_{\XX_{\Gamma_\bullet},\rho})$。\tn{底本此处把 $h_0,h_1,h_2$ 全排成下标，与印刷页 405 所用的上标记号不一致；译文照录原文（印刷页 406）。}

在 §8 与 §9 中，我们在 $\Gamma$ 是（至少密切关联于）局部或整体绝对伽罗瓦群时应用这些想法。（把它们推广到投射有限群的情形在 [Zhu25a, §2.4] 中处理。）为此我们回忆：局部伽罗瓦上同调在次数 $>2$ 时总是消没；整体伽罗瓦上同调也是如此，只要我们在奇剩余特征中工作（我们将始终如此）。

\section*{参考文献}

\begin{refs}
\item[{[AC14]}] Patrick B. Allen and Frank Calegari, \emph{Finiteness of unramified deformation rings}, Algebra Number Theory \textbf{8} (2014), no. 9, 2263--2272, DOI 10.2140/ant.2014.8.2263. MR3294389
\item[{[AGH19]}] Benjamin Antieau, David Gepner, and Jeremiah Heller, \emph{$K$-theoretic obstructions to bounded $t$-structures}, Invent. Math. \textbf{216} (2019), no. 1, 241--300, DOI 10.1007/s00222-018-00847-0. MR3935042
\item[{[AHR20]}] Jarod Alper, Jack Hall, and David Rydh, \emph{A Luna \'etale slice theorem for algebraic stacks}, Ann. of Math. (2) \textbf{191} (2020), no. 3, 675--738, DOI 10.4007/annals.2020.191.3.1. MR4088350
\item[{[All19]}] Patrick B. Allen, \emph{On automorphic points in polarized deformation rings}, Amer. J. Math. \textbf{141} (2019), no. 1, 119--167, DOI 10.1353/ajm.2019.0003. MR3904768
\item[{[Alp13]}] Jarod Alper, \emph{Good moduli spaces for Artin stacks} (English, with English and French summaries), Ann. Inst. Fourier (Grenoble) \textbf{63} (2013), no. 6, 2349--2402, DOI 10.5802/aif.2833. MR3237451
\item[{[Alp14]}] Jarod Alper, \emph{Adequate moduli spaces and geometrically reductive group schemes}, Algebr. Geom. \textbf{1} (2014), no. 4, 489--531, DOI 10.14231/AG-2014-022. MR3272912
\item[{[ALBM24]}] Johannes Ansch\"utz, Arthur-C\'esar Le Bras, and Lucas Mann, \emph{A 6-functor formalism for solid quasi-coherent sheaves on the Fargues--Fontaine curve}, arXiv:2412.20968, 2024.
\item[{[And21]}] Grigory Andreychev, \emph{Pseudocoherent and perfect complexes and vector bundles on analytic adic spaces}, arXiv:2105.12591.
\item[{[Ant21]}] Benjamin Antieau, \emph{On the uniqueness of infinity-categorical enhancements of triangulated categories}, arXiv:1812.01526.
\item[{[Art89]}] James Arthur, \emph{Unipotent automorphic representations: conjectures: Orbites unipotentes et repr\'esentations, II}, Ast\'erisque \textbf{171-172} (1989), 13--71. MR1021499
\item[{[ATJLPRVG18]}] Leovigildo Alonso Tarr\'io, Ana Jerem\'ias L\'opez, Marta P\'erez Rodr\'iguez, and Mar\'ia J. Vale Gonsalves, \emph{On the derived category of quasi-coherent sheaves on an Adams geometric stack}, J. Pure Appl. Algebra \textbf{222} (2018), no. 4, 828--845, DOI 10.1016/j.jpaa.2017.05.009. MR3720855
\item[{[BB10]}] Laurent Berger and Christophe Breuil, \emph{Sur quelques repr\'esentations potentiellement cristallines de $\mathrm{GL}_2(\Q_p)$} (French, with English and French summaries), Ast\'erisque \textbf{330} (2010), 155--211. MR2642406
\item[{[BC08]}] Laurent Berger and Pierre Colmez, \emph{Familles de repr\'esentations de de Rham et monodromie $p$-adique} (French, with English and French summaries), Repr\'esentations $p$-adiques de groupes $p$-adiques. I. Repr\'esentations galoisiennes et $(\phi,\Gamma)$-modules, Ast\'erisque \textbf{319} (2008), 303--337. MR2493221
% 参考文献续于 ch-a4-21.tex（自 [BC09] 起）；refs 环境跨片，由 ch-a4-21 闭合。
